\textbf{适用条件：} 椭圆、双曲线中心原点，焦点横轴；抛物线顶点原点、轴为横轴。抛物线仅关于自身轴对称。
\textbf{公式：} 椭圆 $x^2/a^2+y^2/b^2=1$，a>c>0，b²=a²-c²；双曲线 $x^2/a^2-y^2/b^2=1$，c>a>0，b²=c²-a²；焦点均(±c,0)。抛物线 $y^2=\pm2px$，p>0，焦点(±p/2,0)，准线x=∓p/2。纵向交换坐标。
\textbf{简证：} 两焦半径平方差4cx；椭圆已知和2a，双曲线右支已知差2a，由和差求一个距离后只平方一次；再从方程反向取正根或绝对值验证。抛物线非负距离一次等价平方。
\textbf{例与边界：} F=(±3,0)，距离和10给椭圆 $x^2/25+y^2/16=1$；距离差4给双曲线 $x^2/4-y^2/5=1$。距离差8大于焦距6，轨迹为空，不能写成双曲线。a为椭圆长半轴或双曲线实半轴，c为半焦距。
